% Generated draft: requires mathematical and source review.
\documentclass[11pt]{article}
\usepackage[margin=25mm]{geometry}
\usepackage{fontspec}
\setmainfont{Latin Modern Roman}
\usepackage{amsmath,amssymb,mathtools,unicode-math}
\setmathfont{Latin Modern Math}
\usepackage{enumitem,needspace,xurl,hyperref}
\setlength{\parindent}{0pt}
\setlength{\parskip}{5pt}
\setlength{\emergencystretch}{3em}
\newcommand{\N}{\mathbb{N}}
\newcommand{\Z}{\mathbb{Z}}
\newcommand{\Q}{\mathbb{Q}}
\newcommand{\R}{\mathbb{R}}
\newcommand{\C}{\mathbb{C}}
\newcommand{\F}{\mathbb{F}}
\newcommand{\A}{\mathbb{A}}
\newcommand{\PP}{\mathbb{P}}
\newcommand{\E}{\mathbb{E}}
\newcommand{\eps}{\varepsilon}
\newcommand{\dd}{\,\mathrm{d}}
\begin{document}
\section*{UnsolvedMath ID: 30003230}
\textbf{Source problem code:} catalogue entry
\subsection{Definitions and mathematical statement}
Write \(e_{\mathrm{alg}}^{\mathrm{str}}(Y)\) for the algebraic stringy Euler number of a variety \(Y\). The supplied record names this invariant but does not specify its definition, normalization, or hypotheses for its existence and finiteness.

In standard Mori-theoretic terminology, a divisorial Mori contraction is a projective birational contraction of a \(K_X\)-negative extremal ray whose exceptional locus is a divisor. This terminology presupposes an appropriate setting for the canonical divisor and the Mori cone; the record does not specify the ground field, the singularity assumptions, or the projectivity and \(\mathbb Q\)-factoriality hypotheses on the varieties.

The proposed assertion is that, for every divisorial Mori contraction \(f\colon X\to X'\) in the intended setting, with both displayed invariants defined and finite,
\[
 e_{\mathrm{alg}}^{\mathrm{str}}(X)
 >
 e_{\mathrm{alg}}^{\mathrm{str}}(X').
\]
The question is whether this strict inequality always holds. The intended setting and the definition of the invariant require confirmation from the supplied references before this becomes a fully specified mathematical statement.

The supplied dated curation labels the problem partially solved and reports an established flip inequality in the equivariant Mori program for projective spherical varieties, while describing the unrestricted conjecture as open. That assessment concerns flips explicitly, whereas the displayed question concerns divisorial contractions; its precise relevance to this question is unresolved here.
\subsection{Short English statement}
Does every divisorial Mori contraction strictly decrease the algebraic stringy Euler number? The source leaves the precise class of varieties and the definition and finiteness conditions for the invariant unspecified. Its supplied partial-solution assessment discusses flips, so its application to divisorial contractions needs review.
\subsection{Sources}
\begin{itemize}
\item Local UnsolvedMath registry, supplied complete record: permanent ID \(30003230\), source problem code \texttt{catalogue entry}, title \emph{Stringy Euler Numbers under Divisorial Contractions}. The record gives proposed year 2016 and supplies no named proposer.
\item Supplied report citation: \emph{Singularities (2017)}, Oberwolfach Reports, pp.\ 2625--2691, DOI \texttt{10.4171/owr/2016/46}. \url{https://doi.org/10.4171/owr/2016/46}. Also listed in the supplied literature sources as \emph{Oberwolfach Report 46}.
\item Batyrev--Gagliardi, \emph{On the algebraic stringy Euler number}, supplied literature reference: \url{https://arxiv.org/abs/1610.03842}.
\item Literature assessment checked on \texttt{2026-08-21}: the supplied curation labels the problem partially solved, reports the flip inequality for the equivariant Mori program of projective spherical varieties, and describes the unrestricted conjecture as open. This is supplied dated curation derived from the Oberwolfach material, not an independently verified literature assessment; its time-sensitive status and cited claims require independent verification.
\end{itemize}
\end{document}
